\documentclass[conference]{IEEEtran}
\IEEEoverridecommandlockouts
% The preceding line is only needed to identify funding in the first footnote. If that is unneeded, please comment it out.
\usepackage{cite}
\usepackage{amsmath,amssymb,amsfonts}
\usepackage{algorithmic}
\usepackage{graphicx}
\usepackage{textcomp}
\usepackage{xcolor}
\def\BibTeX{{\rm B\kern-.05em{\sc i\kern-.025em b}\kern-.08em
    T\kern-.1667em\lower.7ex\hbox{E}\kern-.125emX}}

\begin{document}

\title{Design and Development of Object Recognition and Distance Measurement Device\\}

\author{\IEEEauthorblockN{1\textsuperscript{st} Author Name}
\IEEEauthorblockA{\textit{dept. name of organization (of Aff.)} \\
\textit{name of organization (of Aff.)}\\
City, Country \\
email address or ORCID}
\and
\IEEEauthorblockN{2\textsuperscript{nd} Author Name}
\IEEEauthorblockA{\textit{dept. name of organization (of Aff.)} \\
\textit{name of organization (of Aff.)}\\
City, Country \\
email address or ORCID}
}

\maketitle

\begin{abstract}
Visual impairment and situational awareness challenges affect millions of people worldwide, limiting their ability to perceive and interact with their environment independently. This project presents an intelligent assistive device that combines real-time object recognition with distance measurement to provide auditory feedback about surrounding objects. The system utilizes a Raspberry Pi 4B as the central processing unit, integrated with a Raspberry Pi camera module for live video capture and an ultrasonic sensor for precise distance measurement. The camera feed is processed using the YOLO (You Only Look Once) deep learning algorithm trained on the COCO (Common Objects in Context) dataset, enabling real-time identification of 80 common object classes. When objects are detected, their names are converted to speech and delivered through earphones, providing immediate auditory feedback to the user. Simultaneously, the ultrasonic sensor measures the distance to objects and speaks out the measured value, giving users critical spatial awareness. This dual-modality approach—visual recognition combined with distance measurement—creates a comprehensive environmental awareness system suitable for assisting visually impaired individuals, enhancing situational awareness for professionals in low-visibility environments, and serving as an educational platform for computer vision and embedded AI applications. The system demonstrates how affordable, readily available components can be leveraged to create sophisticated assistive technology with potential to significantly improve quality of life and safety for users.
\end{abstract}

\begin{IEEEkeywords}
Object Recognition, Distance Measurement, Raspberry Pi, YOLO, Assistive Technology, Ultrasonic Sensor
\end{IEEEkeywords}

\section{Introduction}
The ability to perceive and understand one's environment is fundamental to independent living and safe navigation. For visually impaired individuals, this capability is severely limited, creating challenges in mobility, safety, and quality of life. According to the World Health Organization, approximately 2.2 billion people worldwide have some form of vision impairment, with 90\% of the visually impaired population living in developing nations where access to assistive technology is limited. Traditional assistive devices such as white canes and guide dogs provide valuable but limited information about the environment, offering tactile awareness only of immediately adjacent obstacles without identifying what those obstacles are or providing distance information.

This project addresses these limitations by developing an intelligent assistive device that combines computer vision with distance sensing to provide comprehensive environmental awareness through auditory feedback. The system leverages the computational capabilities of the Raspberry Pi 4B—a credit-card-sized computer powerful enough to run state-of-the-art deep learning models while maintaining low power consumption and portability. A Raspberry Pi camera module captures live video feed, which is processed using the YOLOv8 object detection algorithm trained on the COCO dataset, enabling recognition of 80 common object classes including people, vehicles, furniture, and everyday items.

The addition of an ultrasonic distance sensor provides critical spatial information that pure visual recognition cannot offer. While object detection identifies what is present, the ultrasonic sensor measures how far away it is—information essential for safe navigation and interaction. When objects are detected, their names are converted to speech using text-to-speech technology and delivered through earphones. Similarly, distance measurements are spoken aloud, providing users with both identification and localization information.

The integration of these technologies creates a powerful assistive tool with applications extending beyond visual impairment assistance. The system can enhance situational awareness for professionals working in low-visibility environments such as firefighters, search and rescue personnel, or security personnel. It also serves as an excellent educational platform for teaching computer vision, embedded AI, and assistive technology concepts.

This project represents a convergence of embedded systems, deep learning, computer vision, and human-computer interaction. By demonstrating that sophisticated assistive technology can be built with affordable, readily available components, this work contributes to the broader goal of making advanced AI-powered assistance accessible to those who need it most.

\section{Literature Survey}
Several researchers and practitioners have explored the integration of computer vision and spatial sensing for assistive technology and robotics. The literature demonstrates a transition from traditional obstacle detection mechanisms toward AI-driven semantic understanding of environments.

\subsection{Vision-Based Object Recognition on Embedded Systems}
Recent advancements in lightweight deep learning models have made on-device inference feasible. For example, Kristisswaran (2025) demonstrated animal recognition utilizing YOLO and OpenCV on Raspberry Pi 5 with GPU acceleration, validating the capabilities of edge AI \cite{b3}. Similarly, Jaryd (2024) detailed the setup and optimization of YOLO models on Raspberry Pi platforms specifically targeted at real-time recognition \cite{b8}. These works highlight the critical balance between inference speed and detection accuracy, noting that while smaller models achieve usable frame rates, they still face limitations in dynamic environments. In autonomous navigation contexts, wujiafa9 (2024) showcased ``RaspberryCar,'' which successfully fused TensorFlow Object Detection API outputs with ultrasonic range finding for robotics applications \cite{b2}.

\subsection{Assistive Technologies and Text-to-Speech (TTS)}
Providing intuitive feedback to users is a cornerstone of assistive tool design. The OTON GLASS project by the Raspberry Pi Foundation (2018) highlighted the utility of converting captured image data (text via OCR) into synthesized speech for the visually impaired \cite{b4}. Furthering this approach for generalized reading, recent studies in 2025 introduced regional language reading aids on Raspberry Pi that convert text arrays into audible formats with high accuracy \cite{b9}.

\subsection{Distance Estimation and Sensor Fusion}
While deep learning models can estimate distance through monocular vision (as demonstrated by Jyothi et al. in vehicle speed estimation \cite{b10} and Kumar \cite{b5} using MobileNet-SSD), these approaches heavily rely on known physical assumptions and environmental lighting. Acoustic sensors remain popular and robust alternatives. Hazlina (2024) explored proximity-triggered AI object detection \cite{b7}, establishing that relying on external hardware like ultrasonic modules is computationally cheaper and less error-prone in direct line-of-sight than visual estimation. High-precision alternatives are also emerging, such as Time-of-Flight (ToF) cameras compatible with embedded platforms that provide millimeter-accuracy distance matrices \cite{b6}.

\subsection{Summary of Related Existing Solutions}
A tabular summary comparing the most relevant foundational applications combining Raspberry Pi platforms with object detection and measuring logic is outlined in Table \ref{tab1}.

\begin{table}[htbp]
\caption{Summary of Related Existing Solutions}
\begin{center}
\begin{tabular}{|p{0.5cm}|p{2cm}|p{2cm}|p{2.5cm}|}
\hline
\textbf{Ref.} & \textbf{Topic} & \textbf{Author/Year} & \textbf{Highlights} \\
\hline
\cite{b3} & Animal Recognition using OpenCV on RPi 5 & Kristisswaran (2025) & YOLO, OpenCV, COCO dataset utilizing GPU acceleration. \\
\hline
\cite{b5} & Object Detection \& Distance Estimation & R. Kumar (2024) & MobileNet-SSD with distance est. on object heights. \\
\hline
\cite{b8} & YOLO Recognition on Raspberry Pi & Jaryd (2024) & Setup of YOLO vision models on RPi 5. \\
\hline
\cite{b4} & OTON GLASS & RPi Foundation (2018) & Wearable TTS assisting individuals with visual impairment. \\
\hline
\cite{b9} & Regional Reading Aid for Blind People & IEEE Conference (2025) & TTS via RPi achieving 93\% text-to-voice accuracy. \\
\hline
\cite{b10} & Vehicle Detection \& Speed Estimation & B. Jyothi et al. (2025) & YOLOv8 vehicle detection. \\
\hline
\cite{b1} & YOLOv8 Drone using Raspberry Pi 4B & a7m-1st (2024) & Drone-based integration with real-time data transmission. \\
\hline
\cite{b2} & RaspberryCar & wujiafa9 (2024) & Intelligent car integrating ultrasonic and TensorFlow. \\
\hline
\cite{b7} & Project Ideas for Raspberry Pi & Hazlina (2024) & AI object detection using camera and ultrasonic sensors. \\
\hline
\cite{b6} & Time of Flight Camera for Raspberry Pi & Arducam (2024) & ToF up to 4-meters with sub-2cm tracking precision. \\
\hline
\end{tabular}
\label{tab1}
\end{center}
\end{table}

\section{Problem Statement}
Visually impaired individuals face significant challenges in independent navigation and environmental awareness. Traditional assistive tools provide limited information: white canes detect only immediate obstacles through contact, offering no identification of what the obstacle is or its distance beyond arm's reach. Guide dogs offer more flexibility but require extensive training, have limited availability, and cannot provide verbal information about detected objects.

Existing electronic assistive devices often suffer from high costs, limited functionality, or impractical form factors. Ultrasonic canes provide distance information but cannot identify objects. Camera-based systems that attempt object recognition often require cloud connectivity, introducing latency and privacy concerns. Commercial solutions with integrated object recognition and distance measurement typically cost thousands of dollars, placing them beyond reach for many potential users.

Furthermore, professionals working in low-visibility environments—firefighters in smoke-filled buildings, search and rescue teams in disaster zones, or security personnel in darkness—face similar challenges in maintaining situational awareness. There is a critical need for an affordable, portable, real-time system that can identify objects and measure their distance, providing this information through an intuitive auditory interface.

\section{Existing System}
Current assistive technologies for environmental awareness span several categories. White canes remain the most common tool, providing tactile feedback through contact with obstacles. Electronic Travel Aids (ETAs) such as the UltraCane or K-Sonar add ultrasonic distance sensing but provide only obstacle detection without identification.

Smartphone-based applications like Seeing AI or TapTapSee use computer vision to identify objects but require users to actively point and capture images, providing information only on demand rather than continuous awareness. These apps typically rely on cloud processing, introducing latency and requiring internet connectivity.

Research systems have demonstrated various approaches. The OTON GLASS project uses a Raspberry Pi with camera to capture text and convert it to speech, demonstrating the feasibility of portable vision-based assistive technology. Similar systems have been developed for object recognition using YOLO on Raspberry Pi, achieving real-time performance with optimized models.

However, existing systems typically lack integration of both object recognition and distance measurement in a single portable device. Systems that combine both modalities often rely on expensive LiDAR sensors or stereo camera setups that increase cost and computational requirements. There remains a gap in affordable, integrated solutions that provide both object identification and distance information through an intuitive auditory interface.

\section{Proposed System}
The proposed Object Recognition and Distance Measurement Device integrates advanced computer vision with precise distance sensing in a portable, affordable package. The system centers on a Raspberry Pi 4B, which provides sufficient computational power for real-time deep learning inference while maintaining low power consumption suitable for portable operation.

\subsection{Object Recognition Subsystem}
A Raspberry Pi camera module captures continuous live video feed. The video stream is processed using the YOLOv8 (You Only Look Once version 8) object detection algorithm, a state-of-the-art deep learning model optimized for real-time performance. The model is pre-trained on the COCO (Common Objects in Context) dataset, which includes 80 common object categories such as people, vehicles, animals, furniture, and everyday items. The model runs locally on the Raspberry Pi, eliminating the need for internet connectivity or cloud processing and ensuring privacy and low latency.

\subsection{Distance Measurement Subsystem}
An HC-SR04 ultrasonic sensor provides precise distance measurements to objects in the device's field of view. The sensor emits ultrasonic pulses and measures the time taken for echoes to return, calculating distance with millimeter accuracy. The sensor is positioned to align with the camera's field of view, ensuring that distance measurements correspond to the objects being detected visually.

\subsection{Auditory Feedback Subsystem}
A text-to-speech engine running on the Raspberry Pi converts detected object names and measured distances into spoken words. These audio outputs are delivered through standard earphones, providing clear, immediate feedback to the user without disturbing others. The system can be configured to speak object names when first detected, repeat at intervals, or provide continuous distance updates.

\subsection{System Integration}
A Python-based software framework coordinates all subsystems. The camera continuously captures frames, which are passed to the YOLO model for inference. Detected objects are filtered based on confidence thresholds to minimize false positives. The ultrasonic sensor is triggered periodically, with measurements synchronized to the camera's field of view. The text-to-speech engine queues audio outputs, prioritizing recent detections and distance measurements.

\subsection{User Interface}
The system is designed for simple, intuitive operation. Users wear the device (mounted on a harness or integrated into a wearable frame) and earphones. No manual interaction is required during normal operation—the system continuously provides environmental awareness through auditory feedback. Optional controls allow users to adjust speaking rate, enable/disable specific object categories, or request immediate distance measurements.

\section{Objectives}
The primary objectives of the system are as follows:
\begin{itemize}
    \item To design and implement a real-time object recognition system using Raspberry Pi 4B and camera module with YOLOv8 deep learning model trained on COCO dataset.
    \item To integrate ultrasonic distance measurement for providing spatial information about detected objects.
    \item To develop text-to-speech functionality that audibly announces recognized object names and measured distances through earphones.
    \item To create a portable, battery-powered device suitable for continuous field operation.
    \item To optimize system performance for real-time operation on embedded hardware while maintaining detection accuracy.
\end{itemize}

\section{Design Methodology and Working}
\subsection{Hardware Platform Selection}
The Raspberry Pi 4B was selected as the core processor due to its balance of computational power, power efficiency, and ecosystem support. With a quad-core 1.5GHz processor and up to 8GB RAM, it provides sufficient performance for running YOLO models at usable frame rates while maintaining portability. The Raspberry Pi Camera Module V3 provides high-quality image capture with support for the CSI interface, ensuring low-latency video streaming.

\subsection{Software Environment Setup}
The system runs Raspberry Pi OS (64-bit) with a dedicated Python virtual environment to manage dependencies without affecting system stability. The Ultralytics YOLO package is installed, which includes the YOLOv8 model and necessary dependencies including OpenCV for image processing. The virtual environment ensures consistent behavior and simplifies deployment.

\subsection{Object Detection Implementation}
The YOLOv8 nano (\textit{n}) model is selected as the optimal balance between accuracy and speed for embedded deployment. This lightweight version of YOLOv8 achieves approximately 1.5-2 frames per second on Raspberry Pi 4B while maintaining good detection accuracy. The model is pre-trained on the COCO dataset, enabling recognition of 80 common object classes without additional training.

The detection pipeline follows this sequence:
\begin{enumerate}
    \item Camera captures frame (640x480 resolution for performance).
    \item Frame is passed to YOLO model for inference.
    \item Model returns bounding boxes, class labels, and confidence scores for detected objects.
    \item Results are filtered to remove low-confidence detections (threshold typically 0.5).
    \item Detected object names are queued for text-to-speech conversion.
\end{enumerate}

\subsection{Distance Measurement Integration}
The HC-SR04 ultrasonic sensor is interfaced with the Raspberry Pi GPIO pins. A Python library controls the sensor trigger and echo timing. Distance is calculated using the formula:
\begin{equation}
    Distance = \frac{Speed \times Time}{2}
\end{equation}
where $Speed$ is the speed of sound traveling in the air (approx. $343 \text{ m/s}$).

\subsection{Power Management}
The system is designed for battery operation using a portable power bank (minimum 5V/3A output). Power consumption is approximately 5-7 watts during active operation, providing several hours of runtime with a typical 10,000mAh battery pack.

\section{Applications}
\begin{itemize}
    \item \textbf{Visual Impairment Assistance:} Provides real-time environmental awareness for blind and visually impaired individuals, enabling safer navigation and improved independence.
    \item \textbf{Situational Awareness for First Responders:} Assists firefighters, search and rescue personnel, and emergency workers operating in smoke-filled or dark environments.
    \item \textbf{Security and Surveillance:} Enhances perimeter monitoring with automated object detection and distance tracking.
    \item \textbf{Elderly Care:} Helps elderly individuals with declining vision maintain awareness of their surroundings.
    \item \textbf{Robotics Education:} Serves as a platform for teaching computer vision, embedded AI, and sensor integration.
    \item \textbf{Industrial Safety:} Alerts workers to approaching vehicles or equipment in hazardous environments.
    \item \textbf{Museum and Exhibition Guides:} Provides audio descriptions of exhibits as visitors approach.
    \item \textbf{Autonomous Vehicle Research:} Low-cost platform for testing perception algorithms for small-scale autonomous vehicles.
    \item \textbf{Wildlife Monitoring:} Detects and announces animal presence for nature observation or property protection.
    \item \textbf{Assistive Technology Research:} Platform for developing and testing new human-computer interaction paradigms for assistive devices.
\end{itemize}

\section{Advantages and Disadvantages}
\subsection{Advantages}
\begin{itemize}
    \item \textbf{Dual-Modality Sensing:} Combines object recognition (what) with distance measurement (how far) for comprehensive environmental awareness.
    \item \textbf{Real-Time Operation:} Local processing on Raspberry Pi eliminates cloud latency and ensures immediate feedback.
    \item \textbf{Affordable Implementation:} Uses readily available components totaling under \$100 USD.
    \item \textbf{Portable Design:} Battery-powered and compact enough for wearable applications.
    \item \textbf{Privacy-Preserving:} All processing occurs locally; no images are transmitted to external servers.
    \item \textbf{Expandable Architecture:} Additional sensors (thermal camera, GPS, IMU) can be integrated.
    \item \textbf{Intuitive Auditory Interface:} Earphone delivery provides private, continuous feedback.
    \item \textbf{COCO Dataset Coverage:} Recognizes 80 common object classes covering most everyday situations.
    \item \textbf{Open-Source Software:} Based on freely available libraries and models.
    \item \textbf{Educational Value:} Demonstrates practical applications of deep learning and embedded systems.
\end{itemize}

\subsection{Disadvantages}
\begin{itemize}
    \item \textbf{Limited Frame Rate:} Nano models achieve only 1.5-2 FPS on Raspberry Pi 4B, insufficient for fast-moving objects.
    \item \textbf{Battery Life:} Continuous operation limited to 3-5 hours depending on battery capacity.
    \item \textbf{Ultrasonic Sensor Limitations:} Affected by soft surfaces (sound absorption), angled surfaces (deflection), and multiple reflections.
    \item \textbf{Field of View Mismatch:} Camera and ultrasonic sensor have different detection cones, potentially causing misalignment.
    \item \textbf{Processing Power Constraints:} Cannot run larger, more accurate YOLO models in real time locally.
    \item \textbf{Lighting Sensitivity:} Object detection accuracy decreases in low-light conditions.
    \item \textbf{False Positives/Negatives:} Deep learning models may occasionally misidentify objects or miss detections.
    \item \textbf{Audio Feedback Overload:} Continuous speaking may become overwhelming in cluttered environments.
    \item \textbf{No Object Tracking:} System detects objects frame-by-frame without maintaining identity across frames.
    \item \textbf{Learning Curve:} Users need time to interpret and trust the auditory information.
\end{itemize}

\section{System Requirements}
\subsection{Functional Requirements}
\begin{itemize}
    \item Camera shall capture live video at minimum 640x480 resolution for object detection.
    \item YOLOv8 model shall detect objects from COCO dataset (80 classes) with minimum 50\% confidence threshold.
    \item Ultrasonic sensor shall measure distances from 2cm to 400cm with $\pm$3cm accuracy.
    \item System shall process video frames at minimum 1 frame per second for real-time operation.
    \item Text-to-speech engine shall announce detected object names and measured distances within 2 seconds of detection.
    \item Audio output shall be delivered through standard 3.5mm earphone jack.
    \item System shall operate on battery power for minimum 3 hours continuous use.
    \item Ultrasonic sensor shall trigger at minimum 1 Hz frequency for distance updates.
    \item System shall filter out low-confidence detections to minimize false announcements.
    \item User shall be able to adjust speech rate and volume through configuration file.
\end{itemize}

\subsection{Non-Functional Requirements}
\begin{itemize}
    \item System shall operate in temperature range $0^{\circ}$C to $40^{\circ}$C.
    \item All components shall be housed in a portable enclosure suitable for wearable use.
    \item Power consumption shall not exceed 10W during active operation.
    \item System weight including battery shall not exceed 500 grams.
    \item Object detection accuracy shall maintain minimum 70\% for well-lit indoor environments.
    \item Ultrasonic sensor shall function reliably at angles up to 15 degrees from perpendicular.
    \item Software shall be modular to facilitate addition of new sensors or models.
    \item System shall boot and begin operation within 60 seconds of power-on.
    \item Audio announcements shall be clearly intelligible in moderately noisy environments.
    \item System shall include low-battery warning through a distinctive audio tone.
\end{itemize}

\subsection{Hardware Required}
\begin{itemize}
    \item \textbf{Core Processing}: Raspberry Pi 4B, MicroSD Card (32GB minimum, Class 10), 5V/3A Power Bank with appropriate USB cable.
    \item \textbf{Camera}: Raspberry Pi Camera Module V3 (or V2), CSI Camera Cable.
    \item \textbf{Sensing}: HC-SR04 Ultrasonic Distance Sensor.
    \item \textbf{Audio}: Standard 3.5mm earphones/headphones.
\end{itemize}

\subsection{Software Required}
\begin{itemize}
    \item Raspberry Pi OS (64-bit) Bullseye or later.
    \item Python 3.9 or higher.
    \item Relevant Modules: OpenCV, Ultralytics, TTS libraries etc.
\end{itemize}

\section{Future Scope}
The system can be enhanced by integrating a Time-of-Flight (ToF) camera for more precise distance measurement across the entire field of view rather than single-point ultrasonic sensing. Adding thermal camera capability would enable operation in complete darkness and through smoke. Implementing object tracking algorithms would maintain object identity across frames, enabling more natural language descriptions such as ``the person is approaching from the left.'' Integrating GPS would enable location-aware announcements. Adding gesture recognition would allow users to control system parameters through simple hand movements. Developing a bone conduction headphone interface would preserve environmental hearing while providing audio feedback. Implementing scene description generation using large language models would provide higher-level environmental understanding. Optimizing with Raspberry Pi 5 would significantly improve frame rates and enable larger, more accurate models. Adding cellular connectivity would enable remote monitoring for caregivers. Developing machine learning models specifically trained on assistive technology use cases would improve relevance for target users.

\section{Conclusion}
The Object Recognition and Distance Measurement Device developed in this project successfully demonstrates the integration of advanced computer vision with precise distance sensing in an affordable, portable platform. By combining YOLOv8 deep learning on Raspberry Pi 4B with ultrasonic distance measurement and text-to-speech output, the system provides comprehensive environmental awareness through intuitive auditory feedback.

Key achievements include successful deployment of state-of-the-art object detection on embedded hardware, integration of distance sensing synchronized with visual recognition, and development of a practical auditory interface delivering both identification and localization information. The system's ability to recognize 80 common object classes from the COCO dataset while simultaneously measuring distance represents a significant advancement over existing assistive devices that typically provide only one type of information.

The project demonstrates that sophisticated AI-powered assistive technology no longer requires expensive specialized hardware. With total component cost under \$100 USD, this system makes advanced environmental awareness accessible to populations that have traditionally been excluded from technological solutions due to cost barriers. The local processing approach ensures privacy and eliminates dependence on internet connectivity.

While limitations exist—including modest frame rates, ultrasonic sensor constraints, and the need for further optimization—the system provides a solid foundation for continued development. Future enhancements incorporating ToF cameras, thermal imaging, or more powerful processors such as Raspberry Pi 5 could significantly expand capabilities.

This project contributes to the growing field of assistive technology by demonstrating a practical, replicable design that can be adapted and improved by researchers, makers, and organizations serving visually impaired communities. Ultimately, such systems have the potential to dramatically improve independence, safety, and quality of life for millions of individuals with visual impairments worldwide.

\begin{thebibliography}{00}
\bibitem{b1} a7m-1st. (2024). Yolov8-Image-Recognition-Drone-using-Raspberry-pi-4b. GitHub Repository. https://github.com/a7m-1st/Yolov8-Image-Recognition-Drone-using-Raspberry-pi-4b 
\bibitem{b2} wujiafa9. (2024). RaspberryCar: Intelligent car based on Raspberry Pi. GitHub Repository. https://github.com/wujiafa9/RaspberryCar- 
\bibitem{b3} Kristisswaran. (2025). Animal Recognition using OpenCV on Raspberry Pi 5. Cytron.io Tutorial. https://www.cytron.io/tutorial/animal-recognition-using-opencv-on-raspberry-pi-5 
\bibitem{b4} Raspberry Pi Foundation. (2018). OTON GLASS: turning text to speech. Raspberry Pi News. https://www.raspberrypi.com/news/oton-glass/ 
\bibitem{b5} Kumar, R. (2024). Real-Time Object Detection and Distance Estimation using MobileNet-SSD on Raspberry Pi. GitHub Repository. https://github.com/RanjithKumar17904/Real-Time-Object-Detection-and-Distance-Estimation-using-MobileNet-SSD-on-Raspberry-Pi 
\bibitem{b6} Arducam. (2024). Time of Flight Camera for Raspberry Pi. Pimoroni Wholesale. https://wholesale.pimoroni.com/products/time-of-flight-camera-for-raspberry-pi 
\bibitem{b7} Hazlina. (2024). BTE1522 DRE2213 – Week 10 – Project. Blog @ UMP. http://blog.ump.edu.my/hazlina/2024/12/09/bte1522-dre2213-week-10-project/ 
\bibitem{b8} Jaryd. (2024). Getting Started with YOLO Object and Animal Recognition on the Raspberry Pi. Core Electronics Tutorial. https://core-electronics.com.au/guides/raspberry-pi/getting-started-with-yolo-object-and-animal-recognition-on-the-raspberry-pi/ 
\bibitem{b9} IEEE Conference. (2025). Regional Language Reading Aid for Blind People Using Raspberry Pi. IEEE Xplore. https://ieeexplore.ieee.org/document/10991202 
\bibitem{b10} Jyothi, B., et al. (2025). Real-time vehicle detection and speed estimation system using Raspberry Pi and camera module. Bulletin of Electrical Engineering and Informatics, 14(6). https://www.beei.org/index.php/EEI/article/view/9931 
\end{thebibliography}

\end{document}
